
\textbf{Perplexity-based filtering.} A common approach to pre-training data quality filtering scores documents using a reference language model (typically GPT-2) and removes documents with perplexity above a threshold. ProX \citep{Zhou2024ProX} demonstrated average benchmark improvements of approximately 2\% at 1B--100B token scales across C4, DCLM, and FineWeb when applying perplexity-based filtering. DataMan \citep{Peng2025DataMan} studied filtering threshold effects but did not vary model scale as a treatment variable.

\textbf{Near-duplicate removal.} SoftDedup \citep{He2024SoftDedup} reported 1.77\% few-shot improvement over hard deduplication. The theoretical basis for deduplication affecting in-context learning (ICL) stems from Bayesian accounts of ICL \citep{Xie2021Explanation}, which suggest that the frequency of n-gram patterns in pre-training data influences ICL capability. REWIRE \citep{Nguyen2025REWIRE} applied targeted near-duplicate removal and reported 1.0--2.5 percentage point improvements across 22 tasks at 1--7B scale.

\textbf{Domain mixing.} WebOrganizer \citep{Wettig2025Organize} showed a 2.5 percentage point improvement from combining quality filtering with domain mixing at 1B tokens compared to quality filtering alone.

\textbf{Scale interactions.} Chinchilla scaling laws \citep{Hoffmann2022Training} established that optimal compute allocation varies with model size, but these results pertain to compute-optimal training token counts, not data curation strategy. Na et al. [2024] proposed modular approximations for scalable ablations that bridge small-scale experiments to larger-scale predictions. No prior work reports a controlled factorial study of curation strategy $\times$ model scale interactions in pre-training.

\textbf{The FineWeb dataset.} FineWeb \citep{Penedo2025FineWeb2} is an openly released, large-scale web text corpus (sample-10BT subset used here) derived from Common Crawl with quality-filtered preprocessing. It provides reproducible access to web text without the streaming complexity of Dolma v1.7. The Pythia model suite \citep{Biderman2023Pythia} provides models trained at multiple scales under comparable conditions, making it a natural platform for scaling studies. The experiments here use Pythia-style architectures trained from scratch on FineWeb rather than the original Pythia checkpoints.

\textbf{LM evaluation.} Evaluation used the lm-evaluation-harness framework \citep{Gao2021Harness}.
